% answers.tex —— 2011 年·湖北卷·文科 · 参考答案与详细解析片段
% 由 tools/generate_projects.py 依来源总账自 2011/2011_数学真题_湖北文_试题.md + 2011/2011_数学真题_湖北文_解析正文.md + 2011/2011_数学真题_湖北文_答案要点.md 生成，勿手改（改源 MD 后重跑）

\begin{answer}
\gksolhead{1}{}

\textbf{题目}：

已知 $U=\{1,2,3,4,5,6,7,8\}$，$A=\{1,3,5,7\}$，$B=\{2,4,5\}$，则 $\complement_U(A\cup B)=$\mbox{（　　）}

（A）$\{6,8\}$

（B）$\{5,7\}$

（C）$\{4,6,7\}$

（D）$\{1,3,5,6,8\}$

\Needspace{4\baselineskip}
\textbf{解析}
$A\cup B=\{1,2,3,4,5,7\}$，故 $\complement_U(A\cup B)=\{6,8\}$。选 A。

\gkanswer{答案：A。}

\end{answer}
\begin{answer}
\gksolhead{2}{}

\textbf{题目}：

若向量 $\boldsymbol a=(1,2)$，$\boldsymbol b=(1,-1)$，则 $2\boldsymbol a+\boldsymbol b$ 与 $\boldsymbol a-\boldsymbol b$ 的夹角等于\mbox{（　　）}

（A）$-\dfrac\pi4$

（B）$\dfrac\pi6$

（C）$\dfrac\pi4$

（D）$\dfrac{3\pi}4$

\Needspace{4\baselineskip}
\textbf{解析}
\Needspace{5\baselineskip}
$2\boldsymbol a+\boldsymbol b=(3,3)$，$\boldsymbol a-\boldsymbol b=(0,3)$。设夹角为 $\theta$，
\nopagebreak[4]
\[\cos\theta=\frac{9}{3\sqrt2\cdot3}=\frac{\sqrt2}2,\qquad\theta=\frac\pi4.\]
选 C。

\gkanswer{答案：C。}

\end{answer}
\begin{answer}
\gksolhead{3}{}

\textbf{题目}：

若定义在 $\mathbf R$ 上的偶函数 $f(x)$ 和奇函数 $g(x)$ 满足 $f(x)+g(x)=\mathrm e^x$，则 $g(x)=$\mbox{（　　）}

（A）$\mathrm e^x-\mathrm e^{-x}$

（B）$\dfrac12(\mathrm e^x+\mathrm e^{-x})$

（C）$\dfrac12(\mathrm e^{-x}-\mathrm e^x)$

（D）$\dfrac12(\mathrm e^x-\mathrm e^{-x})$

\Needspace{4\baselineskip}
\textbf{解析}
\Needspace{5\baselineskip}
将 $x$ 换成 $-x$：$f(-x)+g(-x)=f(x)-g(x)=\mathrm e^{-x}$。与原式相减：
\nopagebreak[4]
\[2g(x)=\mathrm e^x-\mathrm e^{-x}\Longrightarrow g(x)=\frac12(\mathrm e^x-\mathrm e^{-x}).\]
选 D。

\gkanswer{答案：D。}

\end{answer}
\begin{answer}
\gksolhead{4}{}

\textbf{题目}：

将两个顶点在抛物线 $y^2=2px$（$p>0$）上，另一个顶点是此抛物线焦点的正三角形个数记为 $n$，则\mbox{（　　）}

（A）$n=0$

（B）$n=1$

（C）$n=2$

（D）$n\geqslant3$

\Needspace{4\baselineskip}
\textbf{解析}
由对称性，正三角形另两顶点关于 $x$ 轴对称。设焦点 $F\left(\dfrac p2,0\right)$，$P(x,y)$，则 $PQ=2|y|$，$PF=PQ$。联立 $y^2=2px$ 可得 $y>0$ 与 $y<0$ 两组解，故 $n=2$。选 C。

\gkanswer{答案：C。}

\end{answer}
\begin{answer}
\gksolhead{5}{}

\textbf{题目}：

有一个容量为 $200$ 的样本，其频率分布直方图如图所示，根据样本的频率分布直方图估计，样本数据落在区间 $[10,12)$ 内的频数为\mbox{（　　）}

（A）$18$

（B）$36$

（C）$54$

（D）$72$
\nopagebreak[4]
\begin{center}% 2011 年湖北卷文科第 5 题：容量 200 的样本频率分布直方图。
% 来源 PDF：普通高考/2011/2011湖北文.pdf，SHA-256 97bda477ec8baea4f3f31a9b6013a77e9a4bcd14f462013754d52c02ff5d1770。
% 原页：第 1 页第 1 栏中部，书页码 669。各组距为 2。
\begin{tikzpicture}[>=stealth, semithick, font=\small, x=0.62cm, y=3.2cm]
  \draw[->] (0,0) -- (13,0) node[right] {样本数据};
  \draw[->] (0,0) -- (0,0.22) node[above] {$\dfrac{\text{频率}}{\text{组距}}$};
  \foreach \x/\h in {1/0.02,3/0.02,5/0.05,7/0.15,9/0.19,11/0.09}{
    \draw (\x-1,0) rectangle (\x+1,\h);
  }
  \foreach \x in {0,2,4,6,8,10,12} \node[below] at (\x,0) {\x};
  \foreach \y in {0.02,0.05,0.15,0.19} \node[left] at (0,\y) {\y};
  \node[below left] at (0,0) {$O$};
\end{tikzpicture}
\end{center}
\nopagebreak[4]
\Needspace{4\baselineskip}
\textbf{解析}
区间 $[10,12)$ 对应小矩形的高为 $0.09$（由直方图读出），组距为 $2$，故该组频率为 $0.09\times2=0.18$，频数为 $200\times0.18=36$。选 B。

\gkanswer{答案：B。}

\end{answer}
\begin{answer}
\gksolhead{6}{}

\textbf{题目}：

已知函数 $f(x)=\sqrt3\sin x-\cos x$，$x\in\mathbf R$。若 $f(x)\geqslant1$，则 $x$ 的取值范围为\mbox{（　　）}

（A）$\left\{x\ \middle|\ 2k\pi+\dfrac\pi3\leqslant x\leqslant2k\pi+\pi,\ k\in\mathbf Z\right\}$

（B）$\left\{x\ \middle|\ k\pi+\dfrac\pi3\leqslant x\leqslant k\pi+\pi,\ k\in\mathbf Z\right\}$

（C）$\left\{x\ \middle|\ 2k\pi+\dfrac\pi6\leqslant x\leqslant2k\pi+\dfrac{5\pi}6,\ k\in\mathbf Z\right\}$

（D）$\left\{x\ \middle|\ k\pi+\dfrac\pi6\leqslant x\leqslant k\pi+\dfrac{5\pi}6,\ k\in\mathbf Z\right\}$

\Needspace{4\baselineskip}
\textbf{解析}
\Needspace{5\baselineskip}
$f(x)=2\sin\left(x-\dfrac\pi6\right)\geqslant1$，即 $\sin\left(x-\dfrac\pi6\right)\geqslant\dfrac12$，故
\nopagebreak[4]
\[2k\pi+\frac\pi6\leqslant x-\frac\pi6\leqslant2k\pi+\frac{5\pi}6,\]
即 $2k\pi+\dfrac\pi3\leqslant x\leqslant2k\pi+\pi$。选 A。

\gkanswer{答案：A。}

\end{answer}
\begin{answer}
\gksolhead{7}{}

\textbf{题目}：

设球的体积为 $V_1$，它的内接正方体的体积为 $V_2$，下列说法中最合适的是\mbox{（　　）}

（A）$V_1$ 比 $V_2$ 大约多一半

（B）$V_1$ 比 $V_2$ 大约多两倍半

（C）$V_1$ 比 $V_2$ 大约多一倍

（D）$V_1$ 比 $V_2$ 大约多一倍半

\Needspace{4\baselineskip}
\textbf{解析}
\Needspace{5\baselineskip}
设球半径 $R$，内接正方体棱长 $a$，则 $3a^2=(2R)^2$，$a=\dfrac{2R}{\sqrt3}$。于是
\nopagebreak[4]
\[\frac{V_1}{V_2}=\frac{\frac43\pi R^3}{\frac{8R^3}{3\sqrt3}}=\frac{\pi\sqrt3}{2}\approx2.72.\]
比 $V_2$ 多约 $1.72$ 倍，即约一倍半。选 D。

\gkanswer{答案：D。}

\end{answer}
\begin{answer}
\gksolhead{8}{}

\textbf{题目}：

直线 $2x+y-10=0$ 与不等式组 $\begin{cases}x\geqslant0,\\y\geqslant0,\\x-y\geqslant-2,\\4x+3y\leqslant20\end{cases}$ 表示的平面区域的公共点有\mbox{（　　）}

（A）$0$ 个

（B）$1$ 个

（C）$2$ 个

（D）无数个

\Needspace{4\baselineskip}
\textbf{解析}
可行域顶点为 $(0,0)$，$(0,2)$，$(2,4)$，$(5,0)$。直线 $2x+y=10$ 只经过顶点 $(5,0)$（其余顶点代入分别为 $0,2,8$），故公共点有 $1$ 个。选 B。

\gkanswer{答案：B。}

\end{answer}
\begin{answer}
\gksolhead{9}{}

\textbf{题目}：

《九章算术》“竹九节”问题：现有一根 $9$ 节的竹子，自上而下各节的容积成等差数列，上面 $4$ 节的容积共 $3$ 升，下面 $3$ 节的容积共 $4$ 升，则第 $5$ 节的容积为\mbox{（　　）}

（A）$1$ 升

（B）$\dfrac{67}{66}$ 升

（C）$\dfrac{47}{44}$ 升

（D）$\dfrac{37}{33}$ 升

\Needspace{4\baselineskip}
\textbf{解析}
\Needspace{5\baselineskip}
同“竹九节”问题，设公差 $d$，则由 $4a_1+6d=3$、$3a_1+21d=4$ 解得 $d=\dfrac7{66}$，$a_1=\dfrac{3-6d}4$，故
\nopagebreak[4]
\[a_5=a_1+4d=\frac34+\frac52\cdot\frac7{66}=\frac{67}{66}.\]
选 B。

\gkanswer{答案：B。}

\end{answer}
\begin{answer}
\gksolhead{10}{}

\textbf{题目}：

若实数 $a$，$b$ 满足 $a\geqslant0$，$b\geqslant0$，且 $ab=0$，则称 $a$ 与 $b$ 互补。记 $\varphi(a,b)=\sqrt{a^2+b^2}-a-b$，那么 $\varphi(a,b)=0$ 是 $a$ 与 $b$ 互补的\mbox{（　　）}

（A）必要而不充分条件

（B）充分而不必要条件

（C）充要条件

（D）既不充分也不必要条件

\Needspace{4\baselineskip}
\textbf{解析}
若 $ab=0$，不妨 $b=0$，则 $\varphi=\sqrt{a^2}-a=0$。反之，若 $\varphi=0$，则 $\sqrt{a^2+b^2}=a+b$，平方得 $ab=0$。故为充要条件，选 C。

\gkanswer{答案：C。}

\end{answer}
\begin{answer}
\gksolhead{11}{}

\textbf{题目}：

某市有大型超市 $200$ 家、中型超市 $400$ 家、小型超市 $1400$ 家。为掌握各类超市的营业情况，现按分层抽样方法抽取一个容量为 $100$ 的样本，应抽取中型超市 $\underline{\qquad}$ 家。

\Needspace{4\baselineskip}
\textbf{解析}
超市总数为 $200+400+1400=2000$，中型超市占 $\dfrac{400}{2000}=\dfrac15$，故应抽取 $100\times\dfrac15=20$ 家。

\gkanswer{答案：$20$。}

\end{answer}
\begin{answer}
\gksolhead{12}{}

\textbf{题目}：

$\left(x-\dfrac1{3\sqrt x}\right)^{18}$ 的展开式中含 $x^{15}$ 的项的系数为 $\underline{\qquad}$。（结果用数值表示）

\Needspace{4\baselineskip}
\textbf{解析}
通项 $\mathrm C_{18}^k x^{18-k}\left(-\dfrac1{3\sqrt x}\right)^k=\mathrm C_{18}^k\left(-\dfrac13\right)^k x^{18-\frac{3k}2}$。令 $18-\dfrac{3k}2=15$ 得 $k=2$，系数为 $\mathrm C_{18}^2\cdot\dfrac19=\dfrac{153}9=17$。

\gkanswer{答案：$17$。}

\end{answer}
\begin{answer}
\gksolhead{13}{}

\textbf{题目}：

在 $30$ 瓶饮料中，有 $3$ 瓶已过了保质期。从这 $30$ 瓶饮料中任取 $2$ 瓶，则至少取到 $1$ 瓶已过保质期饮料的概率为 $\underline{\qquad}$。（结果用最简分数表示）

\Needspace{4\baselineskip}
\textbf{解析}
\nopagebreak[4]
\[P=1-\frac{\mathrm C_{27}^2}{\mathrm C_{30}^2}=1-\frac{351}{435}=\frac{28}{145}.\]

\gkanswer{答案：$\dfrac{28}{145}$。}

\end{answer}
\begin{answer}
\gksolhead{14}{}

\textbf{题目}：

过点 $(-1,-2)$ 的直线 $l$ 被圆 $x^2+y^2-2x-2y+1=0$ 截得的弦长为 $\sqrt2$，则直线 $l$ 的斜率为 $\underline{\qquad}$。

\Needspace{4\baselineskip}
\textbf{解析}
\Needspace{5\baselineskip}
圆化为 $(x-1)^2+(y-1)^2=1$，圆心 $(1,1)$、半径 $1$。弦长 $\sqrt2$ 对应圆心到直线距离
\nopagebreak[4]
\[d=\sqrt{1-\left(\frac{\sqrt2}2\right)^2}=\frac{\sqrt2}2.\]
\Needspace{5\baselineskip}
设直线 $y+2=k(x+1)$，即 $kx-y+k-2=0$，则
\nopagebreak[4]
\[\frac{|2k-3|}{\sqrt{k^2+1}}=\frac{\sqrt2}2\Longrightarrow7k^2-24k+17=0,\]
解得 $k=1$ 或 $k=\dfrac{17}7$。

\gkanswer{答案：$1$ 或 $\dfrac{17}{7}$。}

\end{answer}
\begin{answer}
\gksolhead{15}{}

\textbf{题目}：

里氏震级 $M$ 的计算公式为：$M=\lg A-\lg A_0$，其中 $A$ 是测震仪记录的地震曲线的最大振幅，$A_0$ 是相应的标准地震的振幅。假设在一次地震中，测震仪记录的最大振幅是 $1000$，此时标准地震的振幅为 $0.001$，则此次地震的震级为 $\underline{\qquad}$ 级；$9$ 级地震的最大的振幅是 $5$ 级地震最大振幅的 $\underline{\qquad}$ 倍。

\Needspace{4\baselineskip}
\textbf{解析}
\Needspace{5\baselineskip}
$M=\lg1000-\lg0.001=3-(-3)=6$。设 9 级、5 级地震最大振幅分别为 $A_9$、$A_5$，则
\nopagebreak[4]
\[9-5=\lg A_9-\lg A_5=\lg\frac{A_9}{A_5}\Longrightarrow\frac{A_9}{A_5}=10^4=10000.\]

\gkanswer{答案：$6$；$10000$。}

\end{answer}
\begin{answer}
\gksolhead{16}{}

\textbf{题目}：

设 $\triangle ABC$ 的内角 $A$，$B$，$C$ 所对的边分别为 $a$，$b$，$c$。已知 $a=1$，$b=2$，$\cos C=\dfrac14$。

（1） 求 $\triangle ABC$ 的周长；

（2） 求 $\cos(A-C)$ 的值。

\Needspace{4\baselineskip}
\textbf{解析}
\gksubpart{(1)}
$c^2=a^2+b^2-2ab\cos C=1+4-2\cdot1\cdot2\cdot\dfrac14=4$，故 $c=2$，周长 $5$。

\gksubpart{(2)}
\Needspace{5\baselineskip}
$\sin C=\dfrac{\sqrt{15}}4$，$\cos A=\dfrac{b^2+c^2-a^2}{2bc}=\dfrac78$，$\sin A=\dfrac{\sqrt{15}}8$。于是
\nopagebreak[4]
\[\cos(A-C)=\frac78\cdot\frac14+\frac{\sqrt{15}}8\cdot\frac{\sqrt{15}}4=\frac{11}{16}.\]

\gkanswer{答案（1）：$5$。}

\gkanswer{答案（2）：$\dfrac{11}{16}$。}

\end{answer}
\begin{answer}
\gksolhead{17}{}

\textbf{题目}：

成等差数列的三个正数的和等于 $15$，并且这三个数分别加上 $2$，$5$，$13$ 后成为等比数列 $\{b_n\}$ 中的 $b_3$，$b_4$，$b_5$。

（1） 求数列 $\{b_n\}$ 的通项公式；

（2） 数列 $\{b_n\}$ 的前 $n$ 项和为 $S_n$，求证：数列 $\left\{S_n+\dfrac54\right\}$ 是等比数列。

\Needspace{4\baselineskip}
\textbf{解析}
\gksubpart{(1)}
\Needspace{5\baselineskip}
设三个正数为 $5-d$，$5$，$5+d$。分别加 $2$，$5$，$13$ 得 $7-d$，$10$，$18+d$，成等比：
\nopagebreak[4]
\[10^2=(7-d)(18+d)\Longrightarrow d^2+11d-26=0\Longrightarrow d=2\ \text{\text{或}}\ -13.\]
三个数均为正，故 $d=2$，三数为 $3,5,7$，加后为 $5,10,20$。故 $b_3=5$，公比 $2$，$b_n=5\cdot2^{n-3}$。

\gksubpart{(2)}
\Needspace{5\baselineskip}
$b_1=\dfrac54$，$b_2=\dfrac52$，故
\nopagebreak[4]
\[S_n=\frac{\frac54(2^n-1)}{2-1}=\frac54(2^n-1),\]
\Needspace{5\baselineskip}
从而
\nopagebreak[4]
\[S_n+\frac54=\frac54\cdot2^n=5\cdot2^{n-2},\]
是首项 $\dfrac52$、公比 $2$ 的等比数列。

\gkanswer{答案（1）：$b_n=5\cdot2^{n-3}$。}

\gkanswer{答案（2）：证明见解析。}

\end{answer}
\begin{answer}
\gksolhead{18}{}

\textbf{题目}：

如图，已知正三棱柱 $ABC-A_1B_1C_1$ 的底面边长为 $2$，侧棱长为 $3\sqrt2$，点 $E$ 在侧棱 $AA_1$ 上，点 $F$ 在侧棱 $BB_1$ 上，且 $AE=2\sqrt2$，$BF=\sqrt2$。

（1） 求证：$CF\perp C_1E$；

（2） 求二面角 $E-CF-C_1$ 的大小。
\nopagebreak[4]
\begin{center}% 2011 年湖北卷文科第 18 题：正三棱柱 ABC-A1B1C1，E 在 AA1、F 在 BB1。
% 来源 PDF：普通高考/2011/2011湖北文.pdf，SHA-256 97bda477ec8baea4f3f31a9b6013a77e9a4bcd14f462013754d52c02ff5d1770。
% 原页：第 1 页第 3 栏中部，书页码 669。底面边长 2，侧棱 3√2。
\begin{tikzpicture}[>=stealth, semithick, font=\small, scale=0.95]
  \coordinate (A) at (0,0);
  \coordinate (B) at (3.0,0);
  \coordinate (C) at (1.5,1.0);
  \coordinate (A1) at (0,3.0);
  \coordinate (B1) at (3.0,3.0);
  \coordinate (C1) at (1.5,4.0);
  \coordinate (E) at ($(A)!0.667!(A1)$);
  \coordinate (F) at ($(B)!0.333!(B1)$);
  \draw[thick] (A)--(B)--(C)--cycle;
  \draw[thick] (A1)--(B1)--(C1)--cycle;
  \draw[thick] (A)--(A1) (B)--(B1) (C)--(C1);
  \draw (C)--(F);
  \draw (C1)--(E);
  \draw (E)--(F);
  \fill (E) circle (1.2pt) node[left] {$E$};
  \fill (F) circle (1.2pt) node[right] {$F$};
  \node[left] at (A) {$A$};
  \node[right] at (B) {$B$};
  \node[below] at (C) {$C$};
  \node[left] at (A1) {$A_1$};
  \node[right] at (B1) {$B_1$};
  \node[above] at (C1) {$C_1$};
\end{tikzpicture}
\end{center}
\nopagebreak[4]
\Needspace{4\baselineskip}
\textbf{解析}
\Needspace{5\baselineskip}
以 $B$ 为原点，$BC$ 为 $x$ 轴，平面 $ABC$ 为 $xy$ 平面建立坐标系：
\nopagebreak[4]
\[B(0,0,0),\ C(2,0,0),\ A(1,\sqrt3,0),\ B_1(0,0,3\sqrt2),\ C_1(2,0,3\sqrt2),\]
$E(1,\sqrt3,2\sqrt2)$，$F(0,0,\sqrt2)$。

\gksubpart{(1)}
$\overrightarrow{CF}=(1,-\sqrt3,\sqrt2)$，$\overrightarrow{C_1E}=(-1,\sqrt3,-\sqrt2)$，点积 $-1+3-2=0$，故 $CF\perp C_1E$。

\gksubpart{(2)}
\Needspace{5\baselineskip}
平面 $ECF$ 的法向量
\nopagebreak[4]
\[\boldsymbol n_1=\overrightarrow{CE}\times\overrightarrow{CF}=(\sqrt6,3\sqrt2,2\sqrt3),\]
\Needspace{5\baselineskip}
平面 $CFC_1$ 的法向量
\nopagebreak[4]
\[\boldsymbol n_2=\overrightarrow{CF}\times\overrightarrow{CC_1}=(0,-3\sqrt2,0)\sim(0,1,0).\]
\Needspace{5\baselineskip}
故
\nopagebreak[4]
\[\cos\theta=\frac{|\boldsymbol n_1\cdot\boldsymbol n_2|}{|\boldsymbol n_1||\boldsymbol n_2|}
=\frac{3\sqrt2}{6}=\frac{\sqrt2}2,\]
二面角 $E-CF-C_1$ 为 $45^\circ$。

\gkanswer{答案（1）：证明见解析。}

\gkanswer{答案（2）：$45^\circ$。}

\end{answer}
\begin{answer}
\gksolhead{19}{}

\textbf{题目}：

提高过江大桥的车辆通行能力可改善整个城市的交通状况。在一般情况下，大桥上的车流速度 $v$（单位：千米/小时）是车流密度 $x$（单位：辆/千米）的函数。当桥上的车流密度达到 $200$ 辆/千米时，造成堵塞，此时车流速度为 $0$；当车流密度不超过 $20$ 辆/千米时，车流速度为 $60$ 千米/小时。研究表明：当 $20\leqslant x\leqslant200$ 时，车流速度 $v$ 是车流密度 $x$ 的一次函数。

（1） 当 $0\leqslant x\leqslant200$ 时，求函数 $v(x)$ 的表达式；

（2） 当车流密度 $x$ 为多大时，车流量（单位时间内通过桥上某观测点的车辆数，单位：辆/小时）$f(x)=x\cdot v(x)$ 可以达到最大，并求出最大值。（精确到 $1$ 辆/小时）

\Needspace{4\baselineskip}
\textbf{解析}
\gksubpart{(1)}
\Needspace{5\baselineskip}
$0\leqslant x\leqslant20$ 时 $v=60$。$20\leqslant x\leqslant200$ 时由 $v(20)=60$、$v(200)=0$ 得 $v=\dfrac{200-x}3$。故
\nopagebreak[4]
\[v(x)=\begin{cases}60,&0\leqslant x\leqslant20,\\[4pt]\dfrac{200-x}3,&20<x\leqslant200.\end{cases}\]

\gksubpart{(2)}
\Needspace{5\baselineskip}
$0\leqslant x\leqslant20$ 时 $f=60x\leqslant1200$。$20\leqslant x\leqslant200$ 时
\nopagebreak[4]
\[f(x)=\frac{x(200-x)}3\leqslant\frac{10000}3\approx3333,\]
等号在 $x=100$ 时成立。故 $x=100$ 时车流量最大，约 $3333$ 辆/小时。

\gkanswer{答案（1）：$v(x)=\begin{cases}60,&0\leqslant x\leqslant20,\\[4pt]\dfrac{200-x}3,&20<x\leqslant200.\end{cases}$}

\gkanswer{答案（2）：$x=100$ 时最大，约 $3333$ 辆/小时。}

\end{answer}
\begin{answer}
\gksolhead{20}{}

\textbf{题目}：

设函数 $f(x)=x^3+2ax^2+bx+a$，$g(x)=x^2-3x+2$，其中 $x\in\mathbf R$，$a$，$b$ 为常数，已知曲线 $y=f(x)$ 与 $y=g(x)$ 在点 $(2,0)$ 处有相同的切线 $l$。

（1） 求 $a$，$b$ 的值，并写出切线 $l$ 的方程；

（2） 若方程 $f(x)+g(x)=mx$ 有三个互不相同的实根 $0$，$x_1$，$x_2$，其中 $x_1<x_2$，且对任意的 $x\in[x_1,x_2]$，$f(x)+g(x)<m(x-1)$ 恒成立，求实数 $m$ 的取值范围。

\Needspace{4\baselineskip}
\textbf{解析}
\gksubpart{(1)}
\Needspace{5\baselineskip}
$g(2)=0$，$g'(x)=2x-3$，$g'(2)=1$。由公共切线：
\nopagebreak[4]
\[f(2)=8+8a+2b+a=0,\qquad f'(2)=12+8a+b=1.\]
即 $9a+2b=-8$，$8a+b=-11$，解得 $a=-2$，$b=5$。切线为 $y=x-2$。

\gksubpart{(2)}
\Needspace{5\baselineskip}
此时 $f(x)+g(x)=x^3-3x^2+2x$，方程化为
\nopagebreak[4]
\[x[x^2-3x+(2-m)]=0.\]
三个不同实根要求判别式 $1+4m>0$ 且 $2-m\neq0$，即 $m>-\dfrac14$。
\Needspace{5\baselineskip}
设 $q(x)=x^2-3x+2-m$，则 $x_1,x_2$ 为 $q$ 的两根，$q(x)\leqslant0$ 在 $[x_1,x_2]$ 上成立。所证不等式等价于
\nopagebreak[4]
\[h(x)=xq(x)+m<0.\]
因 $x_1,x_2>0$（$x_1+x_2=3$，$x_1x_2=2-m>0$），故 $xq(x)\leqslant0$，从而 $h(x)\leqslant m$。只需 $m<0$，此时 $h(x)<0$ 恒成立。
结合 $m>-\dfrac14$，得 $m\in\left(-\dfrac14,0\right)$。

\gkanswer{答案（1）：$a=-2$，$b=5$，切线 $l:y=x-2$。}

\gkanswer{答案（2）：$\left(-\dfrac14,0\right)$。}

\end{answer}
\begin{answer}
\gksolhead{21}{}

\textbf{题目}：

平面内与两定点 $A_1(-a,0)$，$A_2(a,0)$（$a>0$）连线的斜率之积等于非零常数 $m$ 的点的轨迹，加上 $A_1$，$A_2$ 两点所成的曲线 $C$ 可以是圆、椭圆或双曲线。

（1） 求曲线 $C$ 的方程，并讨论 $C$ 的形状与 $m$ 值的关系；

（2） 当 $m=-1$ 时，对应的曲线为 $C_1$；对给定的 $m\in(-1,0)\cup(0,+\infty)$，对应的曲线为 $C_2$。设 $F_1$，$F_2$ 是 $C_2$ 的两个焦点。试问：在 $C_1$ 上，是否存在点 $N$，使得 $\triangle F_1NF_2$ 的面积 $S=|m|a^2$？若存在，求 $\tan\angle F_1NF_2$ 的值；若不存在，请说明理由。

\Needspace{4\baselineskip}
\textbf{解析}
\gksubpart{(1)}
\Needspace{5\baselineskip}
设 $P(x,y)$，$x\neq\pm a$，则
\nopagebreak[4]
\[\frac{y}{x+a}\cdot\frac{y}{x-a}=m\Longrightarrow y^2=m(x^2-a^2),\]
即 $x^2-\dfrac{y^2}{m}=a^2$。
$m>0$ 为双曲线；$-1<m<0$ 为焦点在 $x$ 轴的椭圆；$m=-1$ 为圆 $x^2+y^2=a^2$；$m<-1$ 为焦点在 $y$ 轴的椭圆。

\gksubpart{(2)}
\Needspace{5\baselineskip}
$C_1:x^2+y^2=a^2$。对 $m\in(-1,0)\cup(0,+\infty)$，$C_2$ 的半焦距满足 $c^2=a^2(1+m)$，$F_1F_2=2a\sqrt{1+m}$。设 $N(x_0,y_0)\in C_1$，则
\nopagebreak[4]
\[S=\frac12F_1F_2|y_0|=a\sqrt{1+m}|y_0|=|m|a^2,\]
\Needspace{5\baselineskip}
故 $|y_0|=\dfrac{|m|a}{\sqrt{1+m}}$。存在条件 $|y_0|\leqslant a$ 化为
\nopagebreak[4]
\[m^2\leqslant1+m\Longrightarrow\frac{1-\sqrt5}2\leqslant m\leqslant\frac{1+\sqrt5}2.\]
\Needspace{5\baselineskip}
结合 $m\in(-1,0)\cup(0,+\infty)$，得
\nopagebreak[4]
\[m\in\left[\frac{1-\sqrt5}2,0\right)\cup\left(0,\frac{1+\sqrt5}2\right].\]
\Needspace{5\baselineskip}
又 $\cos\angle F_1NF_2=\dfrac{a^2-c^2}{NF_1\cdot NF_2}$，由面积公式 $NF_1\cdot NF_2=\dfrac{2|m|a^2}{\sin\angle F_1NF_2}$，$a^2-c^2=-ma^2$，故
\nopagebreak[4]
\[\tan\angle F_1NF_2=\frac{2|m|a^2}{-ma^2}.\]
$m<0$ 时为 $2$，$m>0$ 时为 $-2$。

\gkanswer{答案（1）：$C:x^2-\dfrac{y^2}{m}=a^2$（$m\neq0$）；$m>0$ 双曲线，$-1<m<0$ 椭圆，$m=-1$ 圆，$m<-1$ 椭圆。}

\gkanswer{答案（2）：存在条件 $m\in\left[\dfrac{1-\sqrt5}2,0\right)\cup\left(0,\dfrac{1+\sqrt5}2\right]$；$m<0$ 时 $\tan\angle F_1NF_2=2$，$m>0$ 时 $-2$。}

\end{answer}